\documentclass[11pt]{article}
\usepackage[top=0.7in,bottom=0.7in,left=0.9in,right=0.9in]{geometry}
\usepackage[T1]{fontenc}
\usepackage{lmodern,amsmath,amssymb}
\usepackage[hidelinks]{hyperref}
\pagestyle{empty}\setlength{\parindent}{0pt}\setlength{\parskip}{0.6em}
\begin{document}
\hfill November 2, 2026

Editor-in-Chief\\ACM Transactions on Algorithms

Dear Editor-in-Chief,

Please consider the enclosed manuscript, ``Dynamic Betti Numbers Are Equivalent to Dynamic Sparse Matrix Rank'', for publication in \emph{ACM Transactions on Algorithms}.

Persistent homology is increasingly applied to data that change over time, but for general complexes of dimension two or more every known method for updating Betti numbers or barcodes after a single simplex insertion or deletion takes time polynomial in the size of the complex in the worst case. The manuscript explains why. Its main theorem is an equivalence: maintaining the first Betti number of a $2$-complex under simplex insertions and deletions and maintaining the rank of a sparse matrix under entry updates reduce to each other with constant overhead (over $\mathbb{F}_2$, and over every field for matrices with entries in $\{0,\pm1\}$). The reduction rests on a local and switchable gadget, adapted from a static construction of Edelsbrunner and Parsa, in which an entry flip costs seven simplex updates and switching a row or column on or off costs one.

Two kinds of consequences follow.
\begin{itemize}
\item \emph{Lower bounds.} Under the OMv conjecture, no data structure with polynomial preprocessing time maintains any Betti number $\beta_p$ ($p\ge1$) with amortized update time $O(n^{1/2-\varepsilon})$ and query time $O(n^{1-\varepsilon})$. The bound holds over every field. It also holds for persistent Betti numbers and barcodes under function-value updates, for online zigzag persistence, on complexes of bounded degree, on $2$-complexes linearly embedded in $\mathbb{R}^4$, for clique and Vietoris--Rips complexes of changing metric spaces and of changing point sets in $\mathbb{R}^{O(\log n)}$, and even for deciding whether the first homology group vanishes. Under the hinted conjecture of van den Brand, Nanongkai, and Saranurak, it becomes a worst-case bound of $n^{\alpha/2-o(1)}$ over every field, where $\alpha/2\ge5/8$, and $\alpha/2\approx0.703$ for the current bounds on matrix multiplication.
\item \emph{Algorithms.} Every dynamic-rank algorithm transfers to Betti numbers. In addition, all Betti numbers of subcomplexes of any fixed complex linearly embedded in $\mathbb{R}^3$ can be maintained in $O(\log^2 n)$ amortized time per update, through a combinatorial form of Alexander duality and dynamic connectivity, so maintaining Betti numbers is easy in three dimensions, whereas for $2$-complexes in $\mathbb{R}^4$ it is as hard as maintaining the rank of a sparse $0/1$ matrix (and equivalent to it over $\mathbb{F}_2$); this is the dynamic counterpart of the static picture of Edelsbrunner and Parsa (SODA 2014). Without geometry, given a tree decomposition of the Hasse diagram of bounded width, polylogarithmic worst-case updates are possible unconditionally.
\end{itemize}

The work connects dynamic algorithms, fine-grained complexity, and computational topology, and I believe it will interest TALG readers in all three areas. All proofs are complete and appear in the body of the paper. The combinatorial identities used in the proofs were also checked by brute-force computation; the code is included as supplementary material.

A preliminary version is available as an arXiv preprint. The manuscript has not been published, and it is not under consideration at any other venue.

Sincerely,\\[1.2em]
Shawn Ray\\Carnegie Mellon University
\end{document}
